\documentstyle[% 


righttag% 


,11pt% 


,amssymb% 


,russian%               remove in Eng. 


,izvru%                 Set izven% in Eng. 


]{amsart} 


 


% 


 


\newcommand{\col}{\operatorname{col}} 


\newcommand{\tts}[1]{} 


 


%\hoffset=-15mm%                  For Eng. 


\setcounter{page}{3} 


\god{2003} 


\nomer{4~(491)} %                        Remove in Eng. 


%\nomer{Vol.\,47, No.\,4}%               Set in Eng. 


%\str{\pageref{begin}--\pageref{end}}%  Set in Eng. 


\udk{517.929} 


 


\author{\it А.Р.\,Абдуллаев, Н.А.\,Брагина} 


% NB:   ^^ remove in Eng. 


\title{Операторы Грина с минимальной нормой} 


 


\date{} 


%\pagestyle{myheadings}%         Set in Eng. 


\pagestyle{plain} %              Remove in Eng. 


\begin{document} 


%\thispagestyle{plain}%          Set in Eng. 


\maketitle 


\label{begin} 


 


 


В теории функционально-дифференциальных уравнений (ФДУ) возникает  


следующая задача: для данного линейного уравнения среди всех краевых задач  


найти такую, оператор Грина которой имеет минимальную норму. Впервые 


вопрос  


об операторе Грина с минимальной нормой рассматривался в [1] в связи с  


проблемой параметризуемости множества решений квазилинейного уравнения.  


Одно из приложений рассматриваемой задачи имеется в [2], где предложена  


схема исследования на разрешимость краевых задач для квазилинейных ФДУ,  


основанная на использовании вспомогательной линейной краевой задачи. При 


этом  


оказывается, что эффективность получаемых признаков разрешимости зависит 


от 


выбора вспомогательной задачи, точнее, зависит от нормы оператора Грина 


этой 


задачи. Отметим также работу [3] с условиями, при 


которых минимальную норму имеет оператор Грина задачи Коши (оператор 


Коши). 


 


B предлагаемой статье дано полное решение поставленной задачи в случае,  


когда пространство решений уравнения является гильбертовым. 


 


\section{Правый обратный оператор} 


 


Пусть $X$, $Y$, $Z$ --- банаховы пространства, $\Lambda:X\to Y$ --- 


линейный ограниченный  


оператор с ядром $\ker\Lambda$ и образом $R(\Lambda)$. Зафиксируем 


некоторый проектор 


$P:X\to X$ на $\ker\Lambda$, и пусть $P^c=I-P$ --- дополнительный 


проектор. 


 


\begin{defn} 


Пусть $\Lambda:X\to Y$ --- линейный сюръективный оператор.  


Линейный оператор $K_p:Y\to X$, удовлетворяющий условиям 


$$ 


1) \ \Lambda K_p=I:Y\to Y,\qquad 2) \ K_p\Lambda=P^c, 


$$ 


называется правым обратным для $\Lambda$, соответствующим проектору $P$. 


\end{defn} 


 


Отметим, что в определении речь идет об алгебраической обратимости. Если  


ядро $\ker\Lambda$ дополняемо в $X$ и $P:X\to X$ --- ограниченный 


проектор с образом  


$R(P)=\ker\Lambda$, то оператор $K_p$ является ограниченным [4]. Далее, 


говоря о правом 


обратном операторе, будем предполагать его ограниченность.  


 


Рассмотрим разложение 


$$ 


X=\ker\Lambda\oplus X_p,\quad 


X_p=\ker P, 


$$ 


порождаемое фиксированным проектором $P$ на $\ker\Lambda$ и положим  


$\pi_p:X\to X_p$, $\pi_p x=P^c x$. Тогда для оператора $\Lambda$ 


справедливо 


мультипликативное представление 


\begin{equation} 


\Lambda=\Lambda_p\cdot\pi_p, 


\end{equation} 


где $\Lambda_p:X_p\to Y$ --- изоморфизм. 


 


Пусть $i_p:X_p\to X$, $i_p x=x$, --- оператор естественного вложения. 


 


\begin{lem} 


Если $K_p:Y\to X$ --- правый обратный оператор для $\Lambda$, соответствующий 


проектору $P$, то 


$$ 


1) \ K_p=i_p\Lambda^{-1}_p,\qquad 


2) \ \|K_p\|=\|\Lambda^{-1}_p\|. 


$$ 


\end{lem} 


 


{\bf Доказательство.} Справедливость утверждения 1) леммы проверяется 


непосредственно с помощью равенств


$$ 


\pi_p\cdot i_p=I:X_p\to X_p,\quad 


i_p\cdot \pi_p=P^c. 


$$ 


Далее имеем 


$$ 


\|K_p\|=\sup_{\|y\|=1}\|i_p\Lambda^{-1}_p y\|= 


\sup_{\|y\|=1}\|\Lambda^{-1}_p y\|=\|\Lambda^{-1}_p\|. 


\quad\square 


$$ 


 


\section{Оператор Грина абстрактной линейной краевой задачи} 


 


Рассмотрим систему двух линейных уравнений 


\begin{equation} 


L x=y, \quad l x=\alpha, 


\end{equation} 


где $L:X\to Y$ --- линейный ограниченный оператор и $l:X\to R^n$ --- 


линейный 


ограниченный вектор-функционал. Следуя ([5], с.\,103), 


систему (2) будем называть 


абстрактной линейной краевой задачей (АЛКЗ). Известно, что многие классы  


линейных краевых задач для обыкновенных дифференциальных уравнений, а 


также  


для функционально-дифференциальных уравнений допускают запись в виде 


АЛКЗ. 


 


Предположим, что задача (2) однозначно разрешима для любых пар правых  


частей. Тогда имеет место представление решения $x=U\alpha+G_l y$. 


Здесь $U$ --- фундаментальный вектор решений однородного уравнения $L 


x=0$,  


удовлетворяющий условию $l(U)=E$, где $E$ --- единичная матрица порядка 


$n\times n$.  


Оператор $G_l:Y\to X$ ставит в соответствие каждому $y\in Y$ единственное 


решение 


краевой задачи 


$$ 


L x=y,\quad l x=0. 


$$ 


Этот оператор называется [5] {\it оператором Грина} задачи (2). 


 


\begin{lem} 


Оператор $P:X\to X$, $P x=U l x$, является проектором на $\ker L$.  


Оператор Грина $G_l$ является правым обратным для оператора $L$, 


соответствующим проектору $P$. 


\end{lem} 


 


Задачу об операторе Грина с минимальной нормой сформулируем в  


следующем виде. Рассмотрим линейное уравнение 


\begin{equation} 


L x=y 


\end{equation} 


с сюръективным оператором $L$. Множество всех операторов Грина для 


уравнения  


(3) обозначим через $\bold G(L)$. Отметим, что эти операторы порождаются 


только функционалами, которые удовлетворяют условиям [4] 


$$ 


1) \ R(l)=R^n,\qquad 2) \ \ker l\oplus\ker L=X. 


$$ 


 


Положим 


$$ 


\mu(L)=\inf\{\|G\|:G\in\bold G(L)\}. 


$$ 


Если для некоторого оператора $G\in\bold G(L)$ выполняется равенство 


$\mu(L)=\|G\|$, то  


будем говорить, что он имеет минимальную норму. 


 


\section{Коэффициент сюръективности} 


 


\begin{defn} 


Коэффициентом сюръективности (КС) оператора $L:X\to Y$ 


называется неотрицательное число, определенное равенством 


$$ 


q(L)=\inf_{z\ne0}\frac{\|L^*z\|}{\|z\|}, 


$$ 


где $L^*:Y^*\to X^*$ --- сопряженный оператор.  


Отметим, что неравенство $q(L)>0$ эквивалентно сюръективности оператора. 


\end{defn} 


 


Через $\cal L(X,Y)$ обозначим пространство линейных ограниченных 


операторов. Основные свойства КС содержит 


 


\begin{lem} 


Пусть $L,T\in\cal L(X,Y)$, $S\in\cal L(Y,Z)$ и $\lambda\in R$, тогда 


справедливы следующие утверждения${:}$ 


 


$1)$ $0\le q(L)\le\|L\|;$ 


 


$2)$ $q(\lambda L)=|\lambda|q(L);$ 


 


$3)$ $q(L+T)\le q(L)+\|T\|;$ 


 


$4)$ если $L$ --- изоморфизм, то $q(L)=\|L^{-1}\|^{-1};$ 


 


$5)$ $q(S)q(L)\le q(S L)\le q(S)\|L\|$. 


\end{lem} 


 


\begin{thm} 


Для коэффициента сюръективности справедлива оценка 


\begin{equation} 


\frac{1}{q(L)}\le\|K_p\|. 


\end{equation} 


Если $\|P^c\|=1$, то $\|K_p\|=\frac{1}{q(L)}$. 


\end{thm} 


 


\begin{pf} 


Так как $K_p$ --- правый обратный оператор для $L$, то $L K_p=I$. 


Поэтому $1=q(I)=q(L K_p)\le q(L)\|K_p\|$. 


Отсюда получаем неравенство (4). 


 


Покажем справедливость второй части утверждения леммы. Применяя  


свойство 5) коэффициента сюръективности к представлению (1), получаем 


\begin{equation} 


q(L)\le q(L_p)\|\pi_p\|. 


\end{equation} 


Так как $L_p$ --- изоморфизм, то 


\begin{equation} 


q(L_p)=\|L^{-1}_p\|^{-1}=\|K_p\|^{-1}. 


\end{equation} 


Здесь мы воспользовались свойством 4) коэффициента сюръективности и  


утверждением 2) леммы 1. По предположению $\|P^c\|=1$, поэтому 


$\|\pi\|=\|P^c\|=1$. 


Теперь из (5) и (6) получаем неравенство $\|K_p\|\le1/q(L)$. 


Это неравенство в сочетании с (4) дает требуемый результат для 


правого обратного оператора. 


\end{pf} 


 


Таким образом, минимально возможным значением нормы правого  


обратного оператора для данного линейного оператора является величина 


$(q(L))^{-1}$. 


Если построен такой проектор $P$ на $\ker L$, что $\|P^c\|=1$, то 


соответствующий правый  


обратный имеет минимальную норму. 


 


\section{Существование оператора Грина с минимальной нормой} 


 


Пусть $D$ --- гильбертово пространство со скалярным произведением 


$(\cdot,\cdot)$, $B$ --- 


банахово пространство и $D$ изоморфно прямому произведению $B\times R^n$, 


$L:D\to B$ --- 


линейный ограниченный оператор, $l:D\to R^n$ --- линейный ограниченный  


вектор-функционал. 


 


Рассмотрим задачу (2) в предположении однозначной разрешимости ее для 


любых  


пар правых частей. Справедлива 


 


\begin{thm} 


Пусть $U=\{x_1,\dotsc,x_n\}$ --- ортонормированный базис $\ker L$ и 


вектор-функ\-ци\-о\-нал  


определен равенством $l x=\col\{l_1x,l_2x,\dotsc,l_n x\}$, $l_i 


x=(x_i,x)$. Тогда 


оператор Грина краевой задачи $(2)$ имеет минимальную норму. 


\end{thm} 


 


{\bf Доказательство.} Рассмотрим проектор $P:X\to X$, определенный  


равенством $P x=U l x=\sum\limits^n_{i=1}x_i(x_i,x)$. 


Этот проектор является ортогональным, т.\,к. выполняется равенство 


\begin{multline*} 


(P x,x-P x)=\bigg(\sum^n_{i=1}x_i(x_i,x),x\bigg)-\bigg( 


\sum^n_{i=1}x_i(x_i,x),\sum^n_{j=1}x_j(x_j,x)\bigg)= \\ 


% 


=\sum^n_{i=1}(x_i,x)^2-\sum^n_{i=1} 


(x_i,x_i)(x_i,x)^2=\sum^n_{i=1}(x_i,x)^2- 


\sum^n_{i=1}(x_i,x)^2=0. 


\end{multline*} 


 


Оператор Грина задачи (2) совпадает с правым обратным для $L$,  


соответствующим проектору~$P$. Так как для ортогонального проектора 


$\|P\|=\|P^c\|=1$, то можно применить теорему 1, из которой и следует 


требуемое утверждение. 


 


Таким образом, в соответствии с утверждением теоремы 2 для построения 


требуемой краевой задачи для данного уравнения достаточно описания базиса 


однородного уравнения. В этом смысле это утверждение достаточно 


эффективно.  


Кроме того, теорема 2 дает принципиальную возможность вычисления  


минимальной нормы оператора Грина без его построения: для этого 


достаточно найти $(q(L))^{-1}$. 


 


Следующее вспомогательное утверждение может оказаться полезным в 


случае, если $B$ --- гильбертово пространство. 


 


Пусть $B=H$ --- гильбертово пространство, 


$[\delta,r]:D\to H\times R^n$, $[\delta,r]x=\{\delta x,r x\}$, --- 


фиксированный изоморфизм. Справедлива  


 


\begin{lem} 


Пространство $D$ является гильбертовым пространством  


относительно скалярного произведения 


$$ 


(x,y)_D\overset{\mathrm{def}}{=}{}(\delta x,\delta y)_B+(r x,r y)_{R^n}. 


$$ 


\end{lem} 


 


Пусть $L_2=L_2[0,1]$ --- действительное гильбертово пространство 


суммируемых с квадратом функций со скалярным произведением  


$(x,y)_{L_2}=\int\limits^1_0 x(s)y(s)d s$. Тогда $D=D_2[0,1]$ --- 


пространство таких абсолютно 


непрерывных функций, что $\Dot x\in L_2[0,1]$, является гильбертовым 


пространством 


относительно скалярного произведения 


$(x,y)_D=x(0)y(0)+\int\limits^1_0\Dot x(s)\Dot y(s)d s$. 


Пространство $W_2=W_2[0,1]$ абсолютно непрерывных функций таких, что  


$\Ddot x\in L_2[0,1]$, является гильбертовым пространством относительно 


скалярного 


произведения $(x,y)_{W_2}=x(0)y(0)+\Dot x(0)\Dot y(0)+\int\limits^1_0 


\Ddot x(s)\Ddot y(s)d s$. 


Эти пространства рассматриваются ниже. 


 


\begin{exam} 


Рассмотрим сингулярное дифференциальное уравнение второго 


порядка $\Ddot x(t)-\frac{1}{t}\Dot x(t)=f(t)$ с решениями в $W_2[0,1]$. 


В качестве ортонормированного 


базиса выбираем $\{1,\frac{t^2}{2}\}$, тогда $l_1x=(1,x)=x(0)$, 


$l_2x=(\frac{t^2}{2},x)=\int\limits^1_0\Ddot x(t)d t=\Dot x(1)-\Dot x(0)$. 


 


Согласно утверждению теоремы 2 минимальную норму имеет оператор  


Грина краевой задачи 


$$ 


L x=\Ddot x(t)-\frac{1}{t}\Dot x(t),\quad 


l_1x=x(0),\quad l_2(x)=\Dot x(1)-\Dot x(0). 


$$ 


\end{exam} 


 


\begin{exam} 


Рассмотрим уравнение $\Ddot x(t)+x(t)=f(t)$, $x\in W_2[0,\pi]$, и в 


качестве ортонормированного базиса ядра оператора $L x=\Ddot 


x(t)+x(t)=f(t)$ 


выбираем $\left\{\sqrt{\frac{2}{2+\pi}}\sin t,\ 


\sqrt{\frac{2}{2+\pi}}\cos t\right\}$. Поэтому оператор Грина краевой 


задачи 


$$ 


\Ddot x(t)+x(t)=f(t),\quad l_1x=\alpha_1,\quad 


l_2x=\alpha_2, 


$$ 


где $l_1x=\Dot x(0)+\int\limits^\pi_0\Dot x(s)\cos s\,d s$, 


$l_2x=x(0)+\Dot x(0)+\Dot x(\pi)- 


\int\limits^\pi_0 \Dot x(s)\sin s\,d s$, 


имеет минимальную норму. 


\end{exam} 


 


\begin{exam} 


Вопрос об операторе Грина с минимальной нормой для 


уравнения $\Dot x(t)-x(t)=f(t)$ рассматривался в [1]. В 


соответствии с  


утверждением теоремы 2 минимальную норму имеет оператор Грина задачи 


$$ 


\Dot x(t)-x(t)=f(t),\quad 


l x=e x(1)-\int^1_0 e^s x(s)d s. 


$$ 


\end{exam} 


 


\begin{thebibliography}{9} 


 


\bibitem{1} 


Максимов В.П. {\em К вопросу о параметризации множества решений 


функ\-ци\-о\-наль\-но-дифференциальных уравнений} // Функц.-дифференц. 


уравнения. -- Пермь, 1988. -- С.\,14--21. 


\bibitem{2} 


Азбелев Н.В., Максимов В.П. {\em Априорные оценки решений задачи Коши и 


разрешимость краевых задач для уравнений с запаздывающим 


аргументом} // Дифференц. уравнения. -- 1979. -- Т.\,5. 


-- \No\,10. -- С.\,1731--1747. 


\bibitem{3} 


Абдуллаев А.Р. {\em Об операторе Грина с минимальной нормой} // Краевые  


задачи. -- Пермь, 1991. -- C.\,3--6. 


\bibitem{4} 


Абдуллаев А.Р., Бурмистрова А.Б. {\em Элементы теории топологических 


н\"етеровых  


операторов}. -- Челябинск, 1994. -- 93 с. 


\bibitem{5} 


Азбелев Н.В., Максимов В.П., Рахматуллина Л.Ф. {\em Введение в теорию 


функционально-дифференциальных уравнений}. -- М.: Наука, 1991. -- 280 с. 


 


\end{thebibliography} 


 


\vskip 2cm 


 


\post{\it Пермский государственный}{\it Поступила} 


\post{\it технический университет}{30.12.2001} 


 


\label{end} 


\end{document} 


